\begin{onehalfspace}

La prima fase ha mostrato che TOP-4 recupera candidati ancora presenti
nell'istogramma, ma nessun post-processing può ricostruire la coerenza
persa durante l'evoluzione. La seconda fase interviene quindi durante il
calcolo: introduce sindrome e decoder con il codice a ripetizione, passa
allo Steane $[[7,1,3]]$ e raggiunge infine il surface code. La trattazione
formale e i protocolli completi dei primi due laboratori sono raccolti
nella Sezione~\ref{app:laboratori_qec} dell'Appendice~\ref{app:fase2}; qui
restano gli esperimenti, le curve e le conclusioni quantitative.

% -----------------------------------------------
\section{Dalla Ridondanza Ingenua alla Codifica}
\label{sec:ridondanza_codifica}

Il no-cloning impedisce di copiare uno stato ignoto e una misura diretta
ne distruggerebbe la sovrapposizione~\cite{wootters1982}. La \ac{QEC}
codifica invece un qubit logico nelle correlazioni di più qubit fisici e
misura operatori che rivelano l'errore senza rivelare l'informazione
logica~\cite{shorcorrection1995,steane1996,roffe2019}. Questa distinzione
fra copia e codifica è sviluppata nella Sezione~\ref{app:qec_richiami}.

% -----------------------------------------------
\section{Il Ciclo della Correzione: Stabilizzatori, Sindrome, Decoder}
\label{sec:ciclo_qec}

Ogni codice segue il ciclo codifica--rumore--estrazione della
sindrome--decodifica--correzione. Gli stabilizzatori commutano con
l'informazione logica e segnalano l'uscita dal sottospazio valido; il
decoder traduce la sindrome in una correzione fisica o in un aggiornamento
del \emph{Pauli frame}. La distanza $d$ determina quanti errori il codice
può correggere secondo l'Eq.~\ref{eq:distanza_t}; definizioni e
dimostrazione operativa sono nell'Appendice~\ref{app:fase2}.

% -----------------------------------------------
\section{Il Panorama dei Codici e la Scelta Metodologica}
\label{sec:confronto_codici}

La progressione è deliberata: il repetition rende leggibile la sindrome,
lo Steane corregge un errore arbitrario e il surface code introduce soglia
e scalabilità. Il confronto con Shor e five-qubit code è conservato nella
Tabella~\ref{tab:confronto_codici} dell'Appendice~\ref{app:fase2}.

% -----------------------------------------------
\section{Laboratorio I: il Codice a Ripetizione}
\label{sec:repetition}

\subsection{Il codice bit-flip a 3 qubit}

Il codice codifica $\alpha\ket{0}+\beta\ket{1}$ nello stato entangled
$\alpha\ket{000}+\beta\ket{111}$. Le parità $Z_1Z_2$ e $Z_2Z_3$
identificano un singolo errore X senza misurare i dati; il duale nella
base di Hadamard protegge invece da Z. Struttura, circuito di estrazione e
limite X/Z sono descritti nella Sezione~\ref{app:qec_repetition}.

\subsection{Risultati del laboratorio repetition}
\label{subsec:risultati_repetition}

L'iniezione deterministica verifica tutti i casi della lookup table
(Tabella~\ref{tab:sindrome_repetition}). Il test quantitativo usa
$20\,000$ campioni per punto: il codice fallisce quando almeno due dei tre
qubit subiscono un errore, quindi
\begin{equation}
    p_L = 3p^2(1-p) + p^3 = 3p^2 - 2p^3.
    \label{eq:pL_repetition}
\end{equation}

\begin{figure}[H]
\centering
\includegraphics[width=0.75\textwidth]{qec_repetition_curve.pdf}
\caption[Errore logico del codice a ripetizione a 3 qubit]{Errore logico $p_L$ in funzione dell'errore fisico $p$ per il codice a ripetizione a 3 qubit. I punti Monte Carlo del codice bit-flip e del duale phase-flip ($20\,000$ ripetizioni per punto) coincidono con la previsione analitica $3p^2 - 2p^3$ (Eq.~\ref{eq:pL_repetition}) entro l'errore statistico. Nella regione $p < 0.5$ la curva giace sotto la bisettrice $p_L = p$: la codifica rende il qubit logico più affidabile del qubit fisico non protetto. Il punto di \emph{break-even} è a $p = 0.5$, dove correzione e assenza di correzione si equivalgono; oltre, la ridondanza amplifica gli errori anziché sopprimerli.}
\label{fig:qec_repetition}
\end{figure}

L'accordo con l'Eq.~\ref{eq:pL_repetition} è quantitativo su tutto l'intervallo: ad esempio, per $p = 0.10$ si osserva $p_L = 0.0288 \pm 0.0012$ contro il valore teorico $0.0280$, e per $p = 0.20$ si ha $p_L = 0.1026 \pm 0.0021$ contro $0.1040$. Il messaggio operativo è quello che l'intera parte \ac{QEC} svilupperà su codici più potenti: sotto una soglia di errore fisico, la codifica \emph{sopprime} l'errore --- qui da $p$ a $\mathcal{O}(p^2)$ --- mentre sopra la soglia lo \emph{amplifica}. Il codice a ripetizione paga però un limite strutturale già discusso: protegge da un solo tipo di errore per volta. Un errore $Z$ sul codice bit-flip commuta con entrambi gli stabilizzatori, non produce sindrome e resta invisibile --- ragione per cui il passo successivo richiede un codice, come quello di Steane, capace di correggere \emph{qualunque} errore su un qubit.

% -----------------------------------------------
\section{Laboratorio II: il Codice di Steane
\texorpdfstring{$[[7,1,3]]$}{[[7,1,3]]}}
\label{sec:steane}

\subsection{Struttura del codice}

Lo Steane~\cite{steane1996} è un codice \ac{CSS} $[[7,1,3]]$: la
struttura di Hamming separa le sindromi X e Z e consente di correggere un
errore arbitrario su un qubit. Parametri, generatori e lookup table sono
riportati nella Sezione~\ref{app:qec_steane}; nel corpo interessa
verificare che la codifica appartenga al code space e che la soppressione
sia quadratica.

\subsection{Protocollo di laboratorio}

Si preparano $\ket{0_L}$ e $\ket{+_L}$, si iniettano X, Z e Y su tutti i
sette qubit e si stima $p_L$ al crescere di $p$. L'estrazione implementata
non è fault-tolerant: il modello Monte Carlo agisce sui dati e assume la
sindrome ideale. Questa assunzione, insieme al protocollo completo, è
esplicitata nell'Appendice~\ref{app:fase2}.

\subsection{Risultati del laboratorio Steane}
\label{subsec:risultati_steane}

La codifica restituisce sindrome nulla $000000$ su tutti i $2\,000$
campioni, quindi lo stato preparato appartiene al code space. Le sette
sindromi sono distinte e non nulle; la verifica copre tutti i 21 casi X,
Z e Y su sette qubit e conferma la biiettività della
Tabella~\ref{tab:steane_sindrome}.

\paragraph{Il qubit logico batte il fisico: soppressione quadratica.}
Il terzo controllo misura il beneficio della correzione. Sotto un canale depolarizzante di intensità $p$ per qubit, si stima la probabilità di errore logico $p_L$ con una simulazione Monte Carlo nel formalismo di Pauli ($200\,000$ campioni per punto, seed $=42$): per ogni campione si estrae la sindrome, si applica la correzione a peso minimo dalla Tabella~\ref{tab:steane_sindrome} e si verifica se il residuo contiene un operatore logico non banale. La Figura~\ref{fig:qec_steane} mostra il risultato in scala doppio-logaritmica: i punti seguono un andamento $p_L \propto p^2$, con pendenza log-log misurata $1.94$ --- la firma di un codice di distanza $d=3$, che corregge ogni errore singolo e fallisce solo a partire da due errori. Poiché gli eventi di peso due, il cui numero di coppie di qubit scala come $\binom{7}{2}$, forniscono il contributo dominante al fallimento, l'errore logico è soppresso da $\mathcal{O}(p)$ a $\mathcal{O}(p^2)$.

\begin{figure}[H]
\centering
\includegraphics[width=0.72\textwidth]{qec_steane_curve.pdf}
\caption[Errore logico del codice di Steane {$[[7,1,3]]$}]{Errore logico $p_L$ in funzione dell'errore fisico $p$ (canale depolarizzante), scala doppio-logaritmica. I punti Monte Carlo seguono il riferimento $\propto p^2$ (pendenza $d=3$); la bisettrice $p_L = p$ rappresenta il qubit fisico non protetto. L'intersezione individua la \textbf{pseudo-soglia} $p \approx 0.08$: sotto di essa il qubit logico è più affidabile del fisico, sopra la ridondanza amplifica l'errore. Curva rigenerabile da \texttt{figure\_src/gen\_qec\_steane.py}.}
\label{fig:qec_steane}
\end{figure}

Il regime di convenienza è netto: per $p = 0.01$ si ha $p_L = 1.67 \times 10^{-3}$ (già un ordine di grandezza sotto $p$), per $p = 0.05$ si ha $p_L = 3.4 \times 10^{-2} < p$, mentre a $p = 0.10$ si ha $p_L = 0.116 > p$. La pseudo-soglia --- l'intersezione tra la curva e la bisettrice $p_L = p$ --- cade attorno a $p \approx 0.08$: al di sotto, codificare conviene.

\paragraph{Confronto con il repetition e assunzioni.}
Rispetto al repetition code (Sezione~\ref{sec:repetition}), che pure esibisce una soppressione $\propto p^2$, lo Steane paga sette qubit fisici invece di tre ma ottiene un vantaggio qualitativo decisivo: protegge simultaneamente da errori $X$, $Z$ e $Y$, mentre il repetition è cieco a un intero tipo di errore. Va infine ribadita l'assunzione dichiarata nella Sezione~\ref{sec:steane}: l'estrazione di sindrome qui implementata \emph{non è fault-tolerant}. Il modello Monte Carlo inietta errori solo sui qubit dati e assume una lettura di sindrome ideale; un errore sulle ancilla durante l'estrazione può corrompere la sindrome e indurre una correzione errata. La quantificazione di questo effetto --- e le tecniche fault-tolerant che lo mitigano --- appartiene al terreno del surface code, oggetto del Capitolo~\ref{chap:surface}, dove il decoder opera su cicli ripetuti di misura proprio per distinguere gli errori di dato da quelli di misura.

\end{onehalfspace}
